\section{Worum es geht}

\textbf{[SETZUNG]} Dieses Dokument behandelt drei strukturelle ``Hässlichkeiten'' des
Standardmodells und ihre T0-Erklärung: (1) Chiralität der schwachen Wechselwirkung,
(2) Gravitation als separater Tensor-Term, (3) Abwesenheit magnetischer Monopole.
In allen drei Fällen liefert FFGFT eine geometrische Erklärung aus dem Energiefeld
$E_\text{field}(x,t)$ und dem Parameter $\xi$. Methodische Einschränkung: Teile
der Herleitung befinden sich noch in Entwicklung; was als Beweis gilt und was als
Skizze, ist explizit markiert.

\section{Chiralität: Linkshändigkeit der schwachen Wechselwirkung}

\textbf{[B]} Im Standardmodell ist die ausschließliche Linkshändigkeit der schwachen
Wechselwirkung eine Eingabe. In FFGFT hat sie eine geometrische Grundlage, die aus
der Halbzähligkeit auf dem Torus folgt.

\textbf{Halbzähligkeit und Innen/Außen-Asymmetrie.} Spin-$\tfrac12$ folgt aus der
Wicklungszahl $w=n_\phi/n_\theta=\tfrac12$ (Wicklung (1,1) auf $T^4$, 720°-Symmetrie).
Die entscheidende Konsequenz: der Torus hat eine \emph{Innen- und Außenseite} mit
verschiedenen effektiven Kopplungen. Das elektrische Feld zeigt auf der Außenseite
topologisch festgelegt nach innen (Elektron) bzw.\ außen (Positron) -- die Richtung
ist durch die Wicklungsgeometrie vorgegeben, nicht frei wählbar.

\textbf{[B]} Nach der Projektion $T^4\to\mathbb{R}^3\times\mathbb{R}_t$ (A090)
erscheinen Innen- und Außenkopplung als verschiedene chirale Zustände. Die
Projektion ist nicht symmetrisch: die kompakte Topologie wählt eine Kopplungsrichtung
aus. Linkshändige Kopplung ist die Außenperspektive, rechtshändige die
Innenperspektive -- und die schwache Wechselwirkung koppelt ausschließlich an den
projizierten Außenzustand.

\textbf{[S]} \textbf{Was noch nicht ausgeführt ist.} Der qualitative Mechanismus
ist klar: Halbzähligkeit $\Rightarrow$ Innen/Außen-Asymmetrie $\Rightarrow$ chirale
Projektion. Was fehlt: eine explizite Rechnung die zeigt, dass $g_R=0$ (nicht
$g_R=-g_L$) aus dieser Projektion folgt. Das ist die offene quantitative Kante --
nicht die Grundidee.

\section{Gravitation als Feldgradient}

\textbf{[S]} Im SM ist Gravitation ein separates Tensorfeld ($G_{\mu\nu}$). In FFGFT
ist Gravitation der Gradient desselben Energiefelds, das auch Teilchenmassen beschreibt.
Die nichtlineare Feldgleichung:
\[
  \Box E_\text{field} + \xi\,E_\text{field}^3 = 0.
\]

\textbf{[K]} \textbf{Äquivalenz im schwachen Feld.} Für $E_\text{field}=E_0(1+h)$
mit kleiner Störung $h$ linearisiert die Feldgleichung zu einer Form, die mit der
linearisierten Einstein-Gleichung äquivalent ist. Die volle kovariante Tensorformulierung
$g_{\mu\nu}(E_\text{field})$ mit Riemann- und Ricci-Tensor ist im Altkorpus
ausgeführt. Das ist ein \textbf{[K]}-Ergebnis: die Äquivalenz im schwachen Feld ist
gezeigt; die Vollständigkeit der Äquivalenz im starken Feld ist nicht auf
Beweisebene.

\section{Magnetische Monopole}

\textbf{[S]} Magnetische Monopole fehlen im SM; ihre Abwesenheit ist ungeklärt. In FFGFT
sind sie topologische Anregungen des Energiefelds und erfüllen die
Dirac-Quantisierungsbedingung $q_e q_m = 2\pi n\hbar$. Ihre Seltenheit ist eine
Konsequenz der hohen Energieschwelle $\sim E_P/\xi\approx7{,}5\times10^3\,E_P$
-- weit jenseits jeder aktuell erreichbaren Energie.

\textbf{[S]} \textbf{Status.} Das topologische Argument ist plausibel; ein direkter
Beweis, dass das T0-Energiefeld-Monopol exakt die Dirac-Quantisierungsbedingung
erfüllt (und nicht nur reproduziert), ist nicht vollständig ausgeführt.

\section{Einwands-Widerlegungstabelle}

\textbf{[K]} Standardeinwände und ihr Status in FFGFT:

\begin{center}
\renewcommand{\arraystretch}{1.3}
{\small\setlength\tabcolsep{4pt}
\begin{tabular}{p{5.5cm}lp{5cm}}
\toprule
Einwand & Status & Kommentar \\
\midrule
Dimensionsinkonsistenz & Widerlegt & Normierung mit $E_P$, $L_0$ gezeigt \\
Keine GR-Äquivalenz & Teilweise & Schwaches Feld rigoros; starkes Feld Skizze \\
Fehlende Nichtlinearität & Widerlegt & $E^3$-Term vorhanden, Lagrange gezeigt \\
Keine Kovarianz & Widerlegt & Tensor $g_{\mu\nu}$ konstruiert \\
Monopol-Problem & Skizze & Topologische Interpretation, Schwelle plausibel \\
Keine Renormierbarkeit & Widerlegt & $[E^4]=4$, renormierbar \\
\bottomrule
\end{tabular}
}
\renewcommand{\arraystretch}{1.0}
\end{center}

\section{Einordnung in die A-Serie}

\textbf{[B]} Die Chiralitäts-Frage ist in A095 ausgeführt: die Innen/Außen-Topologie
des Torus und die chirale Projektion verbinden die Halbzähligkeit (A263) mit der
Linkshändigkeit der schwachen Kopplung geometrisch.

\textbf{[B]} Die Gravitation-als-Feldgradient-These ist verwandt mit A142 (Zeitfeld-Lagrange):
dort wird Gravitation aus dem Zeitfeld abgeleitet; hier aus dem Energiefeld. Beide
Ansätze müssen auf Konsistenz geprüft werden -- das ist noch nicht geschehen.

\section{Prüfskript}

\begin{itemize}[leftmargin=2em,itemsep=2pt,topsep=2pt]
  \item Chiralitäts-Phasendifferenz für Links-/Rechtshänder; schwache Feld-Linearisierung;
    Monopol-Schwelle aus $E_P/\xi$:
    \texttt{python/A\_Serie\_Skripte/a264\_asymmetrie.py}
\end{itemize}

\section{Was offen bleibt}

\textbf{Offene quantitative Kante: $g_R=0$ aus der Projektion.}
Der Mechanismus Halbzähligkeit $\Rightarrow$ Innen/Außen-Asymmetrie $\Rightarrow$
chirale Projektion ist geometrisch klar. Was noch fehlt: die explizite Rechnung
die zeigt, dass die schwache Kopplung des projizierten Innenzustands genau null
ist -- und nicht nur kleiner als die Außenkopplung. Das ist eine quantitative
Aufgabe, keine konzeptuelle.

\textbf{Gravitation im starken Feld ist nicht auf Beweisebene.} Die schwach-Feld-
Äquivalenz ist gezeigt; die volle nichtlineare Äquivalenz fehlt.

\textbf{Die Verbindung zur Projektionskette (A090) ist nicht ausgeführt.} Ob die
schwache Wechselwirkung aus der Dekompaktifizierung folgt, ist offen.

\section{Ersetzt}

Dieses Dokument nimmt auf: Dok.~143 (Chiralität aus Energiefeld-Rotation, Gravitation
als Feldgradient, magnetische Monopole als topologische Anregungen, Feldgleichung
$\Box E+\xi E^3=0$, kovariant Tensor), Dok.~144 (Einwands-Widerlegungstabelle,
Renormierbarkeit, SM als effektive Niederenergietheorie). \emph{Nicht} übernommen:
ausgeschlossene Anwendungen aus Dok.~143/144 ($\Lambda_\text{cosmo}$-Wert, Spannungsparameter).

\section{Verwendung von KI-Werkzeugen}

Dieses Dokument ist unter Verwendung KI-gestützter Werkzeuge entstanden. Verantwortung
für Inhalt und Fehler liegt beim Autor. Wortlaut der Regeln in A010.
